\documentclass[11pt]{article}
\usepackage[T1]{fontenc}
\usepackage[utf8]{inputenc}
\usepackage{lmodern,microtype,geometry,amsmath,amssymb,amsthm,mathtools,enumitem}
\usepackage[colorlinks=true,linkcolor=blue,urlcolor=blue,citecolor=blue]{hyperref}
\hypersetup{pdftitle={Quadratic-width learning bounds for Bayesian networks with stable neuron populations},pdfauthor={ChatGPT-6 Astra; Claude Fable 5 and Opus 5; Kaarel Hanni; Dmitry Vaintrob}}
\geometry{margin=1in}
\newtheorem{theorem}{Theorem}
\newtheorem{lemma}[theorem]{Lemma}
\newtheorem{proposition}[theorem]{Proposition}
\newtheorem{corollary}[theorem]{Corollary}
\newtheorem{remark}[theorem]{Remark}
\newcommand{\R}{\mathbb R}
\newcommand{\E}{\mathbb E}
\newcommand{\Prb}{\mathbb P}
\newcommand{\norm}[1]{\left\|#1\right\|}
\newcommand{\clip}{\operatorname{clip}}
\title{Quadratic-width learning bounds for Bayesian networks\protect\\with stable neuron populations}
\author{ChatGPT-6 Astra \and Claude Fable 5 and Opus 5 \and Kaarel H\"anni \and Dmitry Vaintrob\thanks{Authors are listed in alphabetical order.}}
\date{Research draft, 21 October 2026}
\begin{document}
\maketitle
\begin{abstract}
We give a sample bound quadratic, up to logarithms, in the width of an
implementing Boolean circuit, for Gaussian Bayesian tanh networks whose
predictions are stable under averaging neurons into redundant populations.
The allowed number of populations is calibrated by the network's readout
energy; the prior restriction does not use the target or its circuit size.
The guarantee holds at every learner width above the teacher width,
including the finite-variance mean-field prior and all stronger scales in
the specified family. For fixed calibration and accuracy, an elementary
bound is $\widetilde O(L^4C^2+L^2d^2)$ examples. A different calibration of
the support condition gives $\widetilde O(L^2C^2+LdC)$.
The proof combines a positive-mass Boolean implementation, Gaussian tails
uniform over all population partitions, and a finite parameter code.
For single-output bounded-fan-in teachers, pruning gives a coarser bound
literally quadratic in circuit width, with a universal degree-nine
polynomial in the other resources.
Population-averaging stability is a strong structural assumption. Whether
it follows from weaker internal-noise robustness, or holds for relevant
trained networks, remains open. This research draft is not Lean verified.
\end{abstract}

\section{Why a different robustness condition helps}

An explicit sample bound with a small polynomial would make a circuit-prior
learning guarantee substantially more informative. Parameter coding suggests
that a width-$C$, depth-$L$ representation should cost roughly
$dC+LC^2$ parameters. The difficulty is that a Bayesian learner can have
arbitrarily larger width $N$: parameter counting for the learner then says
little about the target's complexity.

We study a restriction expressing redundancy among the learner's neurons.
A network must admit averaging into a bounded number of populations while
preserving its clipped predictions. The number of populations is bounded
by readout energy, which the Gaussian posterior can control independently
of learner width. This gives a direct description of each retained posterior
predictor. It bypasses a general entropy-compression argument.

The restriction is stronger and more structural than stability under small
adversarial internal disturbances alone. It therefore defines a new prior
family. We show that every layered Boolean circuit has a positive-mass
implementation in this family, even if the original internal-disturbance
restriction is imposed as well. The theorem concerns ideal Bayesian
conditioning; no efficient posterior inference algorithm is supplied.

\section{Model and population-averaging stability}

Let $d\ge1$, $L\ge2$, and $H=L-1$. Inputs are in
$\mathcal X=\{-1,1\}^d$, with an arbitrary distribution $D$.
All $H$ hidden layers have width $N$ and tanh activation:
\[
 h^1=\tanh(W^1\bar x),\qquad
 h^s=\tanh(W^sh^{s-1}),\qquad
 z_w(x)=v^\top h^H,\qquad \bar x=(x,1).
\]
The raw Gaussian law $\mu_{N,t}$ has independent entries with variances
\[
 \frac{t^2}{d}\quad\hbox{in }W^1,\qquad
 \frac{t^2}{N}\quad\hbox{in }W^s\ (s\ge2),\qquad
 \frac{t^2}{N^2}\quad\hbox{in }v,
 \qquad 0<t\le1.
\]
In particular, $t=N^{(1-\gamma)/2}$ includes every $\gamma\ge1$.
The readout energy and clipped score are
\[
 E_v(w)=N\norm v_2^2,\qquad T(u)=\clip(u,-1/2,1/2).
\]

A balanced partition into $q$ populations has group sizes in
$\{\lfloor N/q\rfloor,\lceil N/q\rceil\}$. Let $P_s$ denote the
orthogonal projection that averages coordinates inside each population at
layer $s$. The population-averaged network $\bar w$ has parameters
\begin{equation}\label{eq:projection}
 \bar W^1=P_1W^1,\qquad
 \bar W^s=P_sW^sP_{s-1}\ (s\ge2),\qquad
 \bar v=P_Hv.
\end{equation}
Each layer is constant inside a population, so this is exactly a width-$q$
tanh network after collecting equal neurons.

Define $\mathcal A_N$ to be the set of networks for which there are an
integer $1\le q\le N$ and balanced population partitions at all hidden
layers satisfying
\begin{equation}\label{eq:population-condition}
 q\le1+E_v(w),\qquad
 \sup_{x\in\mathcal X}|T(z_w(x))-T(z_{\bar w}(x))|\le\frac1{16}.
\end{equation}
The integer $q$ and the partitions can depend on the weights. They cannot
depend on training labels. This is a Borel condition: for fixed $N$ there
are finitely many partitions and finitely many inputs.

We also allow an additional internal-noise stability restriction. Add
disturbances after each hidden activation and at the readout, with norm
\[
 \norm\eta_*^2=\sum_{s=1}^H\frac{\norm{\eta^s}_2^2}{N}
                       +|\eta^{\rm out}|^2,
 \qquad r_c(w)=\frac{c}{\sqrt{1+HE_v(w)}},\quad c>0.
\]
The set $\mathcal R_{N,c}$ consists of networks satisfying
\begin{equation}\label{eq:rms-condition}
 |T(z_w(x;\eta))-T(z_w(x))|\le\frac1{16}
 \quad (x\in\mathcal X,\ \norm\eta_*\le r_c(w)).
\end{equation}
Noisy hidden states need not stay in $[-1,1]^N$.

Let the prior be $\mu_{N,t}$ conditioned on
$\mathcal S_N=\mathcal A_N\cap\mathcal R_{N,c}$.
For a labeled sample $S=((X_i,f(X_i)))_{i=1}^n$, condition further on
\begin{equation}\label{eq:likelihood}
 f(X_i)z_w(X_i)>\frac12\quad(1\le i\le n).
\end{equation}
The resulting posterior is denoted $Q_{S,N,t,c}$.
The same theorem holds with $\mathcal S_N=\mathcal A_N$ alone: the proof
uses \eqref{eq:rms-condition} only to verify that the explicit target
cylinder belongs to this optional smaller support.

\section{An explicit learning theorem}

The target is computed by a depth-$H$, width-$C$ layered Boolean circuit
with NAND, copy, and negation gates. Constant wires needed for internal
gate implementations count toward $C$. Unused stages can contain copies;
the output is one final wire. Compilation from another basis must count
its width and depth overhead. This complexity notion does not refer to
arbitrary real-weight neural circuits.

Define the following target constants:
\begin{align}
 e_*&=\frac1{16\sqrt{2C}}, &\tau_0&=e_*/4,\notag\\
 b_L&=\sqrt{16\log9+8\log(4L)}, &
 A&=\max\{32b_L,2\log(8/e_*)\},\notag\\
 a&=\max\{4,256c\}, &
 \rho&=\min\{1/4,A/[4(d+1)],Ae_*/(64\sqrt C)\},\notag\\
 q_0&=d+2+(H-1)C, &
 B_0&=9dA^2+18(H-1)CA^2+2a^2C,\notag\\
 D_*&=B_0+q_0\log\frac{e\sqrt{2\pi}}{2\rho}+\log2.
 \label{eq:target-constants}
\end{align}
In particular, $D_*\ge1$ and
\begin{equation}\label{eq:target-asymptotic}
 D_*=O\bigl((d+LC)\log^2(2dLC)+c^2C\bigr).
\end{equation}

For $0<\varepsilon,\delta<1$, set
\begin{align}
 B&=\max\{8,8[D_*+\log(4/\varepsilon)]\},
 &Q&=\lceil1+B\rceil,\notag\\
 J&=8\left[D_*+\log(4/\varepsilon)+(H+1)\log Q
       +\frac{\log2}{2}\{d+2+(H-1)Q\}\right],\notag\\
 R&=1+\sqrt{QJ},
 &A_g&=1+d+Q(R+1),\notag\\
 \zeta&=\frac1{32(H+1)A_g^{H+1}},
 &P_Q&=Q(d+2)+(H-1)Q^2.\label{eq:budget}
\end{align}

\begin{theorem}[A quadratic-width sample bound]\label{thm:main}
Suppose $f$ has the Boolean implementation described above. If
\begin{equation}\label{eq:sample}
 n\ge\frac2\varepsilon\left[
 P_Q\log(3+4R/\zeta)+\log(1/\delta)\right],
\end{equation}
then, with probability at least $1-\delta$ over iid inputs $X_i\sim D$,
\[
 \E_{w\sim Q_{S,N,t,c}}\Prb_{X\sim D}
       \{f(X)z_w(X)\le1/4\}\le\varepsilon
\]
simultaneously for every integer $N\ge C$ and every $0<t\le1$.
All posterior normalizers are positive. The conclusion also holds for the
prior restricted only to $\mathcal A_N$.
\end{theorem}

The sample budget has the simpler form
\begin{equation}\label{eq:rate}
 n=\widetilde O\left(
 \frac{LdQ+L^2Q^2+\log(1/\delta)}{\varepsilon}\right),
 \qquad
 Q=\widetilde O\left(d+LC+c^2C+\log(1/\varepsilon)\right).
\end{equation}
The soft-$O$ notation suppresses logarithmic factors in the displayed
resources. For fixed $c$ and fixed accuracy and confidence this is
$\widetilde O(L^2(d+LC)^2)$, and hence quadratic in $C$ up to logarithms.
The dominant depth factor when $LC$ dominates $d$ is $L^4$.
Appendix~\ref{app:calibration} gives a different support calibration with
a smaller depth factor.

There is also a literal quadratic-in-$C$ bound with no hidden logarithms
of $C$, at the cost of a looser polynomial in the other resources. It uses
the bounded fan-in and single-output convention of the Boolean teacher.

\begin{corollary}[A literal quadratic envelope]\label{cor:literal}
There is an absolute constant $K_0$ such that, with
\[
 S_0=2+d+L+c+\log(4/\varepsilon),
\]
the sample condition
\[
 n\ge\frac{K_0}{\varepsilon}
       \left[C^2S_0^9+\log(1/\delta)\right]
\]
suffices for every conclusion of Theorem~\ref{thm:main}, under the same
baseline population condition \eqref{eq:population-condition}.
\end{corollary}
\begin{proof}
Prune the Boolean circuit to the ancestors of its output. At stage $s$,
there are at most $2^{H-s}$ relevant nonconstant wires. Identical constant
wires can be consolidated, leaving at most one additional constant wire
per stage. Thus the same target has an implementation of width
\[
 C'\le\min\{C,2^{H-1}+1\},\qquad \log C'=O(L).
\]
Unused layers retain the prescribed depth by copies. The coarse bound
also covers $H=1$, when no additional hidden constant wire is needed.
Apply the theorem with $C'$ in its target constants; the learner-width
assumption $N\ge C$ remains sufficient.

In \eqref{eq:target-constants}, $A^2=O(S_0^2)$,
$\log(1/\rho)=O(S_0)$, and
$D_*=O(C'S_0^3)$. Consequently
\[
 Q=O(C'S_0^3),\qquad J=O(C'S_0^4),\qquad
 P_Q=O((C')^2S_0^7).
\]
Moreover,
$\log(3+4R/\zeta)=O(L[\log C'+\log S_0])=O(S_0^2)$.
Insert these bounds into \eqref{eq:sample} and use $C'\le C$.
All implicit constants here are absolute.
\end{proof}

\begin{remark}[Prediction conventions]
Let $p_S(x)=Q_{S,N,t,c}\{f(x)z_w(x)>1/4\}$. The conclusion implies
$D\{p_S\le2/3\}\le3\varepsilon$ and
$D\{p_S\le1/2\}\le2\varepsilon$. It also bounds posterior average sign
error. These statements allow arbitrary posterior mixtures.
\end{remark}

\section{A target implementation of sufficient prior mass}

\begin{lemma}[Stable Boolean cylinder]\label{lem:cylinder}
For every $N\ge C$ and $0<t\le1$, there is a set of weights contained in
$\mathcal A_N\cap\mathcal R_{N,c}$ on which $f(x)z_w(x)>1$ for every
input, with raw Gaussian probability at least $\exp(-ND_*/t^2)$.
\end{lemma}
\begin{proof}
Partition the neurons into $C$ balanced clone groups of sizes $k_b$.
Then $C/2\le N/k_b\le2C$. Fill unused logical slots with constants or
copies, so the ideal Boolean state $b^s(x)$ has entries in $\{-1,1\}$,
is constant on each group, and has norm $\sqrt N$. Write $P$ for averaging
within the clone groups.

Use gain $A$ in every gate. NAND has signed affine expression $1-u-v$;
copies and negations use $u$ and $-u$. Count constant wires in internal
layers. Each ideal signed preactivation is at least $A$. Divide each
coefficient from a logical wire equally among its clones, and let the
readout average the output-wire clones with total coefficient $a$.
Denote these teacher weights by $w_*$.

In coordinates $\sqrt d W^1$, $\sqrt N W^s$, and $Nv$, all raw entries
have variance $t^2$. Restrict every first-layer coordinate, every internal
row's coordinate along each normalized clone-group direction, and every
readout coordinate to within $\rho t$ of its teacher value. Internal
matrices then have the form
\[
 W^s=W_*^s+R^s+E^s,\qquad
 R^s=R^sP,\quad \norm{R^s}_{\rm op}\le\rho t\sqrt C,
 \quad E^sP=0.
\]
The unconstrained $E^s$ are independent Gaussian bulk matrices. Two
$1/4$-nets, each of cardinality at most $9^N$, give
\[
 \Prb\{\norm{E^s}_{\rm op}>b_Lt\}
 \le2\exp(2N\log9-Nb_L^2/8)\le\frac1{2L}.
\]
Indeed, each fixed bilinear form has variance at most $t^2/N$, and the
operator norm is at most twice its maximum over the two nets.
All bulk norm bounds therefore hold with conditional probability at least
$1/2$.

We check finite disturbance stability, allowing RMS disturbance $\tau_0$
separately at each hidden layer and scalar disturbance at most $\tau_0$
at the output. This contains the total-budget ball $\norm\eta_*\le\tau_0$.
Put $e_s=\norm{h^s_\eta-b^s}_2/\sqrt N$. If $e_{s-1}\le e_*$, each
group mean differs from its ideal bit by at most
$e_*\sqrt{N/k_b}\le1/16$. At most three teacher coefficients enter a
gate, so its intended signed preactivation remains at least $A/2$.
For $\xi\in\{-1,1\}$ the elementary bound
\[
 \xi q\ge A/2\quad\Longrightarrow\quad
 |\tanh(q+u)-\xi|\le2e^{-A/2}+8|u|/A
\]
follows by considering $|u|<A/4$ and its complement. Since
$E^sb^{s-1}=0$ and $\norm{h^{s-1}_\eta}_2/\sqrt N\le1+\tau_0$,
\[
 e_s\le2e^{-A/2}+(8/A)\rho\sqrt C(1+\tau_0)
                 +(8b_L/A)e_{s-1}+\tau_0\le e_*.
\]
Each term is at most $e_*/4$. In the first layer, the preactivation
perturbation is at most $\rho(d+1)\le A/4$, which initializes the same
invariant. At the readout,
\[
 f(x)z_w(x;\eta)\ge15a/16-\rho(1+\tau_0)-\tau_0>1.
\]
Furthermore, $\sqrt N\norm{v_*}_2\ge a\sqrt{C/2}$ and
$\sqrt N\norm{v-v_*}_2\le\rho$, so
\[
 E_v\ge a^2C/8\ge2C,\qquad
 r_c(w)\le\frac{2\sqrt2c}{a\sqrt C}\le\tau_0.
\]
Thus all noisy and clean clipped outputs on the calibrated ball equal
$f(x)/2$, establishing membership in $\mathcal R_{N,c}$.

The same clone projections in \eqref{eq:projection} map this cylinder
back into itself. Averaging rows preserves the coordinate intervals because
the teacher rows agree within each group. Double projection kills the
bulk, since $E^sP=0$. The first-layer and readout intervals are likewise
preserved. The averaged network therefore also has signed margin greater
than 1 on every input. Its clipped predictions agree exactly with the
original network, and its population count satisfies $C\le1+E_v$.
This proves membership in $\mathcal A_N$.

There are $Nq_0$ independent constrained scalar coordinates; the sum of
squares of their centers is at most $NB_0$. For $Z\sim\mathcal N(0,t^2)$,
$\rho\le1/4$, and any center $m$,
\[
 \Prb\{|Z-m|\le\rho t\}
 \ge\frac{2\rho}{e\sqrt{2\pi}}e^{-m^2/t^2}.
\]
This follows by bounding the Gaussian density throughout the interval,
using $(|m|/t+\rho)^2/2\le m^2/t^2+\rho^2$.
Multiply the interval estimates and the bulk probability $1/2$.
Since $N/t^2\ge1$ and $t\le1$, the result is at least
$\exp(-ND_*/t^2)$ with \eqref{eq:target-constants}.
\end{proof}

\section{Gaussian control uniform over all partitions}

For fixed balanced partitions, let $k_{s,a}$ be the population sizes.
The exactly equivalent width-$q$ network has parameters
\begin{align*}
 U^1_{aj}&=k_{1,a}^{-1}\sum_{i\in a}W^1_{ij},\\
 U^s_{ab}&=k_{s,a}^{-1}\sum_{i\in a,\ j\in b}W^s_{ij},\\
 u_b&=\sum_{j\in b}v_j.
\end{align*}
Define its energy
\[
 \kappa_q=\frac dq\norm{U^1}_F^2+
           \sum_{s=2}^H\norm{U^s}_F^2+q\norm u_2^2,
 \qquad P_q=q(d+2)+(H-1)q^2.
\]

\begin{lemma}[A union bound over adaptive partitions]\label{lem:partitions}
For integers $Q\ge1$ and real $J>0$,
\begin{align*}
 &\mu_{N,t}\{\text{some balanced partitions with }q\le\min(N,Q)
                         \text{ have }\kappa_q>J\}\\
 &\quad\le Q\exp\left(HN\log Q+
      \frac N2[d+2+(H-1)Q]\log2-\frac{NJ}{8t^2}\right).
\end{align*}
\end{lemma}
\begin{proof}
For fixed partitions the compressed coordinates are independent centered
Gaussians, with variances respectively
\[
 \frac{t^2}{dk_{1,a}},\qquad
 \frac{t^2k_{s-1,b}}{Nk_{s,a}},\qquad
 \frac{t^2k_{H,b}}{N^2}.
\]
After applying the factors in $\kappa_q$, each variance is at most
$2t^2/N$. Here balance gives $N/(qk)\le2$, $qk/N\le2$, and a ratio at
most 2 between any two group sizes. Thus
\[
 \E\exp\left(\frac{N\kappa_q}{8t^2}\right)\le2^{P_q/2},\qquad
 \Prb(\kappa_q>J)\le e^{-NJ/(8t^2)}2^{P_q/2}.
\]
There are at most $q^{HN}$ ordered partition tuples. For $q\le N$,
$P_q/N\le d+2+(H-1)q$. Sum over $q\le\min(N,Q)$.
This controls all partitions before selecting one depending on the weights.
\end{proof}

\begin{lemma}[Posterior truncation]\label{lem:truncation}
Under the posterior in Theorem~\ref{thm:main}, outside mass at most
$\varepsilon/2$, there is a witnessing population partition with
$q\le Q$ whose compressed network has energy at most $J$.
\end{lemma}
\begin{proof}
Lemma~\ref{lem:cylinder} bounds the raw conditioning probability below by
$e^{-ND_*/t^2}$, for every correctly labeled sample.
Under the raw Gaussian, $NE_v/t^2$ is chi-square with $N$ degrees of freedom.
For $B\ge8$ its exponential moment gives
\[
 \mu_{N,t}(E_v>B)\le e^{-NB/(4t^2)}2^{N/2}
                     \le e^{-NB/(8t^2)}.
\]
After posterior normalization, the choice of $B$ bounds this probability
by $\varepsilon/4$. If $E_v\le B$, condition
\eqref{eq:population-condition} supplies $q\le Q$.

In Lemma~\ref{lem:partitions}, use $t\le1$ and $N\ge1$ to bound every
positive exponent term by $N/t^2$ times its corresponding term in $J/8$.
The resulting raw probability is at most
\[
 e^{-ND_*/t^2}e^{-(N/t^2)\log(4/\varepsilon)}.
\]
Posterior normalization bounds it by $\varepsilon/4$. Outside these two
exceptional events all partitions of sizes at most $Q$ satisfy the claimed
energy bound, so any witnessing partition can be used.
\end{proof}

\section{Finite coding and generalization}

An energy-$J$ network of width $q\le Q$ has every coefficient bounded by
$\sqrt{QJ}<R$. Pad it to width $Q$ with zero neurons, and round its
parameters to a grid with mesh at most $\zeta$, preserving zero. The rounded
coefficients have absolute value at most $R$. There is a fixed finite class
$\mathcal G$ of all these rounded networks with
\begin{equation}\label{eq:code}
 \log|\mathcal G|\le P_Q\log(3+4R/\zeta).
\end{equation}
For completeness, a grid on $[-R,R]$ with zero included can be chosen with
at most $3+4R/\zeta$ points and rounding error at most $\zeta$.

If $e_s$ is the maximum coordinate error in hidden layer $s$, tanh's
1-Lipschitz property gives
\[
 e_1\le(d+1)\zeta,\qquad e_s\le QR e_{s-1}+Q\zeta.
\]
The output error is at most $QR e_H+Q\zeta$, hence at most
$(H+1)A_g^{H+1}\zeta=1/32$. This is uniform on the Boolean cube.
The original posterior predictor and its rounded candidate therefore
satisfy
\begin{equation}\label{eq:approximation}
 |T(z_w(x))-T(h(x))|\le1/16+1/32=3/32.
\end{equation}

\begin{proof}[Proof of Theorem~\ref{thm:main}]
Define a ramp $\psi$ equal to 1 on $(-\infty,11/32]$, equal to 0 on
$[13/32,\infty)$, and affine between these thresholds. Every retained
posterior network has a candidate $h\in\mathcal G$ satisfying
\eqref{eq:approximation}. Its training losses
$\psi(f(X_i)T(h(X_i)))$ are zero. At any input with
$f(x)z_w(x)\le1/4$, its ramp loss equals 1.

For a fixed bounded nonnegative loss with population mean $r$, the
probability of zero loss on all $n$ iid examples is at most
$(1-r)^n\le e^{-nr}$. Indeed, the probability of a positive loss is at
least its mean. A union bound over $\mathcal G$ shows that, with probability
at least $1-\delta$, every zero-empirical-loss candidate has population
loss at most
\[
 \frac{\log|\mathcal G|+\log(1/\delta)}{n}.
\]
This single sample event is independent of $N$ and $t$. On it,
Lemma~\ref{lem:truncation} and the pointwise loss comparison yield
\[
 \E_Q\Prb_D\{fz_w\le1/4\}
 \le\varepsilon/2+
       \frac{\log|\mathcal G|+\log(1/\delta)}{n}
 \le\varepsilon
\]
simultaneously for every allowed width and scale. All constants defining
$\mathcal G$ depend only on the displayed target and statistical resources.
Finally, \eqref{eq:rate} follows from
$\log(3+4R/\zeta)=O(L\log((d+1)(L+1)(Q+1)(J+1)))$ and
$J=O(d+LQ+L\log Q)$.
\end{proof}

\section{Interpretation and open questions}

The condition says that neuron populations can be replaced by their
averages without substantially changing predictions. It is testable in
principle by searching for suitable neuron clusters, although this paper
provides no efficient certification method. The readout-energy calibration
prevents the restriction from becoming vacuous at arbitrarily large width:
the posterior can afford only a bounded number of retained populations.

The proof supplies an explicit open neighborhood around every cloned
Boolean implementation. It does not establish that all robust low-energy
networks admit such population partitions. Deriving comparable coding bounds
from the original internal-noise condition, or finding a weaker empirically
well-supported condition, is the main unresolved issue.

Quadratic dependence is a useful dense-parameter benchmark, but it need
not be information-theoretically optimal for the teachers considered here.
A bounded-fan-in Boolean circuit has a discrete description of order
$LC\log(C+d)$ bits. The quadratic factor in this theorem comes from coding
a dense real-valued compressed network, not from a lower bound on learning
the Boolean target class.

\paragraph{Acknowledgments and verification status.}
The target-cylinder construction is adapted from the earlier collaborative
work on internally robust Gaussian network priors, including contributions
from Claude. The population-partition Gaussian bound and its application
here are new to this research draft. Independent AI-agent review checked
the main argument; this paper is not Lean verified and has not received a
full human mathematical review.

\appendix
\section{A different calibration with a smaller depth factor}
\label{app:calibration}

Fix a support parameter $\lambda\ge1$ and replace the first inequality of
\eqref{eq:population-condition} by
\begin{equation}\label{eq:lambda-condition}
 q\le1+E_v(w)/\lambda.
\end{equation}
Call this event $\mathcal A_{N,\lambda}$. It is a different prior
restriction: choosing a larger $\lambda$ makes the structural hypothesis
stronger. The parameter is fixed before observing labels.

\begin{corollary}[Calibrated population count]\label{cor:lambda}
Theorem~\ref{thm:main} holds with $\mathcal A_{N,\lambda}$ in place of
$\mathcal A_N$, after making the following changes:
\[
 a=\max\{4\sqrt\lambda,256c\},\qquad
 Q=\lceil1+B/\lambda\rceil.
\]
Recompute $D_*,B,J,R,A_g,\zeta,P_Q$ using these values and the same formulas.
Then the same sample budget \eqref{eq:sample} is sufficient, simultaneously
for all $N\ge C$ and $0<t\le1$. Moreover,
\begin{align*}
 D_*&=O((d+LC)\log^2(2dLC)+(\lambda+c^2)C),\\
 Q&=O\left(1+
 \frac{(d+LC)\log^2(2dLC)+(\lambda+c^2)C+\log(1/\varepsilon)}{\lambda}
 \right).
\end{align*}
\end{corollary}
\begin{proof}
The target-cylinder proof is unchanged except that its readout energy now
satisfies $E_v\ge a^2C/8\ge2\lambda C$, ensuring
$C\le1+E_v/\lambda$. The new readout cost adds $O(\lambda C)$ to $D_*$.
After truncating $E_v\le B$, the support condition gives $q\le Q$ with
the stated smaller $Q$. All Gaussian partition, coding, and statistical
arguments depend on the population cap $Q$, and remain unchanged.
\end{proof}

For the architecture-dependent choice $\lambda=d+L$, fixed $c$, and fixed
accuracy and confidence, the cap is $Q=\widetilde O(C)$, giving
\[
 n=\widetilde O(LdC+L^2C^2).
\]
For arbitrary accuracy, retain
$Q=\widetilde O(C+1+\log(1/\varepsilon)/(d+L))$ in
\eqref{eq:rate}, together with the explicit factors $1/\varepsilon$ and
$\log(1/\delta)$. This sharpening uses a different prior family and is
not a free improvement of the baseline support condition.

\end{document}
